\subsection{Висновки до другого розділу}

У цьому розділі було здійснено перехід від теоретичних математичних моделей до практичної програмної реалізації алгоритму масштабування цифрових зображень. На основі проведеної роботи можна зробити наступні ключові висновки:

\begin{itemize}
	\item \textbf{Узагальнено барицентричний апарат на двовимірний простір.} Математично виведено та обґрунтовано двовимірну форму другої барицентричної формули для просторової сітки Чебишева. Вирішено проблему ділення на нуль при точних збігах цільових координат із вузлами сітки, що гарантує абсолютну обчислювальну стабільність алгоритму.
	
	\item \textbf{Проведено глибоку архітектурну оптимізацію алгоритму.} Доведено практичну неефективність прямого (нативного) методу обчислень через його надмірну асимптотичну складність $\mathcal{O}(W \cdot H \cdot W_2 \cdot H_2)$. Завдяки математичній декомпозиції двовимірної формули та переходу до її матричного подання ($G = B \cdot F \cdot A^{\top}$), обчислювальну складність було кардинально знижено, а сам процес зведено до двох послідовних одновимірних проходів.
	
	\item \textbf{Забезпечено апаратне прискорення обчислень.} Інтеграція бібліотеки OpenCV дозволила делегувати ресурсомісткі матричні множення оптимізованому ядру. Використання векторизації (SIMD-інструкцій), кеш-оптимізації та багатопотоковості перетворило алгоритм на високопродуктивне рішення, здатне обробляти зображення високої роздільної здатності в режимі реального часу.
	
	\item \textbf{Розроблено локальний барицентричний метод.} Для ефективного придушення феномена Гіббса (паразитних осциляцій на різких контурах) реалізовано локальний підхід. Побудова інтерполянтів на ковзних вікнах розміром $K \times K$ пікселів дозволила зберегти структурну цілісність градієнтних переходів і забезпечила гнучкий компроміс між швидкістю та точністю.
	
	\item \textbf{Створено комплексне програмне середовище для тестування.} На базі С++ та фреймворку Qt6 розроблено зручний графічний інтерфейс із підтримкою локальної обробки виділених зон (Region of Interest) та налаштуванням специфічних параметрів (розмір вікна $K$). До системи інтегровано базові алгоритми OpenCV (Nearest, Linear, Cubic, Lanczos4) як конкурентні еталони, що формує надійну експериментальну базу для подальшого об'єктивного порівняння якості масштабування за допомогою формальних метрик.
\end{itemize}